\section{Mitocondri e cloroplasti}

I mitocondri e i cloroplasti sono quegli organuli cellulari in grado di convertire l'energia in forme utilizzabili dalle cellule.

\rubrica{Ciclo dell'energia}
\[
\text{glicolisi} \rightarrow \text{fosforilazione ossidativa} \rightarrow ATP \rightarrow \left\{\text{sintesi proteine, energia meccanica, trasporto attivo, calore}\right\} \rightarrow ADP \rightarrow
\]

\rubrica{Respirazione aerobica e fotosintesi a confronto}
\begin{itemize}
  \item \textbf{respirazione aerobica} (mitocondri, nella maggior parte delle cellule eucariotiche): glucosio $+ O_2 \rightarrow ATP + CO_2 + H_2O$;
  \item \textbf{fotosintesi} (cloroplasti, in alcune cellule vegetali e algali): luce $+ CO_2 + H_2O \rightarrow ATP + O_2 + $ glucosio.
\end{itemize}
I prodotti dell'una sono i substrati dell'altra: il ciclo si chiude tra le due vie.

% ---------------------------------------------------------------------------
\subsection{Mitocondri}

Sono gli organuli respiratori della cellula, deputati alla produzione dell'energia necessaria per molte funzioni cellulari.

\rubrica{Caratteristiche generali}
\begin{itemize}
  \item forma di bastoncino; $\phi$ 0,2\,$\mu$m, lunghezza 2--4\,$\mu$m;
  \item composizione lipoproteica: 65--70\% proteine, 25--30\% lipidi;
  \item presenza di piccole quantità di RNA e DNA;
  \item contengono un elevato contenuto di enzimi capaci di ossidare i prodotti dell'assorbimento intestinale, degradandoli ad $H_2O$ e $CO_2$, liberando energia che viene convertita in energia chimica mediante la fosforilazione di adenosin-di-fosfato (ADP) in adenosin-tri-fosfato (ATP): $ADP + P \rightarrow ATP$.
\end{itemize}

% ---------------------------------------------------------------------------
\subsection{Struttura}

Forma allungata; circoscritti da due membrane:
\begin{itemize}
  \item \textbf{membrana mitocondriale esterna}: decorso regolare, racchiude i mitocondri;
  \item \textbf{membrana mitocondriale interna}: si solleva in numerose pliche denominate \textbf{creste mitocondriali}, orientate trasversalmente all'asse maggiore del mitocondrio.
\end{itemize}
Spazio tra le due membrane $\rightarrow$ \textbf{spazio intermembrana}. Spazio all'interno del mitocondrio $\rightarrow$ \textbf{matrice mitocondriale}.

Sulla membrana interna: particelle di \textbf{ATP sintasi}. Nella matrice: DNA mitocondriale, ribosomi, granuli.

% ---------------------------------------------------------------------------
\subsection{La respirazione mitocondriale}

Le diverse fasi della respirazione aerobica hanno luogo in sedi specifiche:
\begin{itemize}
  \item \textbf{citoplasma}: glicolisi (glucosio $\rightarrow$ piruvato) $\rightarrow$ 2 ATP;
  \item \textbf{ingresso nel mitocondrio}: formazione dell'acetil coenzima A (piruvato $\rightarrow$ acetil-CoA);
  \item \textbf{matrice mitocondriale}: ciclo dell'acido citrico (ciclo di Krebs) $\rightarrow$ 2 ATP;
  \item \textbf{creste mitocondriali}: trasporto degli elettroni e chemiosmosi (catena respiratoria) $\rightarrow$ 32 ATP.
\end{itemize}

% ---------------------------------------------------------------------------
\subsection{Biochimica}

Tre classi di enzimi agiscono in sequenza.

\rubrica{Nella matrice mitocondriale}
\textbf{Enzimi ossidativi del ciclo di Krebs}: degradano i prodotti derivanti dall'assorbimento intestinale ad anidride carbonica ($CO_2$), liberando elettroni (o i loro equivalenti atomi di idrogeno, $H^+$).

\rubrica{Sulla membrana mitocondriale interna}
\begin{itemize}
  \item \textbf{enzimi della catena respiratoria} (sistema di trasporto degli elettroni): trasportano coppie di elettroni (o i loro equivalenti ioni $H^+$), liberati durante il ciclo di Krebs, attraverso varie reazioni intermedie al termine delle quali gli elettroni si combinano con ossigeno molecolare formando acqua ($H_2O$) e cedendo energia;
  \item \textbf{enzimi fosforilativi} (ATP sintasi): catalizzano la sintesi di adenosintrifosfato (ATP) a partire da adenosindifosfato (ADP) e dal fosfato, utilizzando l'energia ceduta da una coppia di elettroni.
\end{itemize}

La catena respiratoria è organizzata in complessi (I, II, III, IV) sulla membrana interna; il gradiente protonico generato attraversa il complesso V (ATP sintasi), producendo ATP a partire da ADP e $P_i$.

% ---------------------------------------------------------------------------
\subsection{Altre funzioni}

\begin{itemize}
  \item metabolismo dei lipidi e dei fosfolipidi;
  \item capacità di ossidare gli acidi grassi;
  \item partecipazione, insieme al reticolo endoplasmatico liscio, alla biosintesi degli ormoni steroidei;
  \item in grado di accumulare e di concentrare varie sostanze;
  \item nella matrice: modesta sintesi proteica e presenza di piccole quantità di DNA, di ribosomi e delle varie classi di RNA.
\end{itemize}

% ---------------------------------------------------------------------------
\subsection{Genoma mitocondriale}

\begin{itemize}
  \item presenza di DNA e di tutto l'apparato per la sintesi proteica: ribosomi, tRNA, mRNA;
  \item DNA mitocondriale simile a quello batterico;
  \item si formano per divisione di mitocondri preesistenti;
  \item si duplicano come i batteri;
  \item le molecole di DNA mitocondriale sono troppo piccole per codificare tutte le proteine di questo organulo $\rightarrow$ i mitocondri sono quindi \textbf{semi-autonomi}, cioè capaci di duplicarsi e presiedere alla sintesi di alcuni costituenti chimici.
\end{itemize}

\rubrica{Analogie con i batteri}
\begin{itemize}
  \item il DNA ha forma circolare e sono assenti le proteine cromosomiche;
  \item i ribosomi sono di grandezza simile o inferiore a quelli dei batteri;
  \item la membrana interna del mitocondrio con il sistema di trasporto degli elettroni è paragonabile alla membrana plasmatica dei batteri, che ha funzione respiratoria e che forma invaginazioni (\textit{mesosomi}) simili alle creste mitocondriali;
  \item sulla base di queste analogie, si ipotizza che i mitocondri possano rappresentare batteri assunti nel corso dell'evoluzione all'interno di cellule primitive e che abbiano realizzato una forma di simbiosi con la cellula ospite.
\end{itemize}

% ---------------------------------------------------------------------------
\subsection{Cloroplasti}

\rubrica{Struttura}
Due membrane:
\begin{itemize}
  \item \textbf{membrana esterna}: porine, $\phi$ 1 nm;
  \item \textbf{membrana interna}: impermeabile (presenza di trasportatori), ripiegata a formare i \textbf{tilacoidi}.
\end{itemize}
\begin{itemize}
  \item i \textbf{tilacoidi} sono disposti in pile ordinate dette \textbf{grana};
  \item spazio interno ai tilacoidi $\rightarrow$ \textbf{lume};
  \item spazio interno del cloroplasto (esterno ai tilacoidi) $\rightarrow$ \textbf{stroma}, che contiene gli enzimi per la sintesi dei carboidrati.
\end{itemize}

Presenza di DNA, residuo di un batterio endosimbionte ancestrale.

% ---------------------------------------------------------------------------
\subsection{Fotosintesi clorofilliana}

\rubrica{Fase luminosa}
Assorbimento dell'energia luminosa, che viene conservata sotto forma di ATP e NADPH. Si ha sviluppo di ossigeno.

\rubrica{Fase oscura}
ATP e NADPH vengono utilizzati per ridurre la $CO_2$ a formare glucosio e altri prodotti organici.

Reazione complessiva:
\[
6CO_2 + 6H_2O \rightarrow C_6H_{12}O_6 + 6O_2.
\]

\rubrica{Fotosistemi}
\begin{itemize}
  \item \textbf{fotosistema I}: complesso proteico (dimero P700) associato all'antenna LHCI; assorbe energia luminosa e, tramite una catena di trasportatori (tra cui la ferredossina), riduce il NADP\textsuperscript{+} a NADPH;
  \item \textbf{fotosistema II}: complesso proteico (dimero P680, proteine D1/D2) associato all'antenna LHCII; il \textbf{complesso sviluppante ossigeno} (con centro Mn--Ca) scinde l'acqua liberando elettroni, protoni e $O_2$;
  \item il trasporto di elettroni tra PSII e PSI avviene attraverso trasportatori mobili (plastochinone PQ, citocromo $b_6f$, plastocianina PC), con generazione di un gradiente protonico attraverso la membrana tilacoidale.
\end{itemize}

% ---------------------------------------------------------------------------
\subsection{Origine di mitocondri e cloroplasti}

Teoria dell'\textbf{endosimbiosi}: cellula ospite di tipo procariotico $\rightarrow$ acquisizione di batteri aerobi (endosimbionti), tramite numerose invaginazioni della membrana plasmatica.
\begin{itemize}
  \item i batteri aerobi diventano \textbf{mitocondri};
  \item il reticolo endoplasmatico e l'involucro nucleare si formano dall'invaginazione della membrana plasmatica (non incluso nell'ipotesi dell'endosimbiosi seriale);
  \item batteri fotosintetici diventano \textbf{cloroplasti}, in cellule eucariotiche vegetali e in alcuni protisti;
  \item le cellule eucariotiche di animali, funghi e alcuni protisti possiedono mitocondri ma non cloroplasti.
\end{itemize}
